%% toc: P2 10 10.1,10.2
\section{Readability and valuations over an arbitrary set of prime formulas}
\label{sec:p02-w02b-p02-001-prime-readability}

Barwise sets up propositional logic over an arbitrary set \(P\) of
\emph{prime formulas} and extends each truth assignment
\(v:P\to\{\mathbf t,\mathbf f\}\) to a function \(\bar v\) on all
propositional formulas of \(P\) ``by induction on length of formulas''
\cite[ch.~A.1, \S4, pp.~23--24]{Barwise1977FirstOrder}; his definitions are
reproduced in \cref{def:p02-w02-p02-005-bw-propformula} and
\cref{def:p02-w02-p02-005-bw-truth-assignment}.  That recursion defines a
function only if every propositional formula of \(P\) is built in exactly
one way, and the chapter neither states nor proves such a property.  The
handbook's readability theorem \cref{thm:prop-unique-readability} does not
settle the point either: it concerns the countable alphabet of
\cref{sec:prop-syntax-semantics-01}, with variables \(p_n\) and the constant
\(\bot\), whereas Barwise allows \(P\) to be any set, for example a set of
sentences of a natural language, and in his application takes \(P\) to be a
set of first-order formulas, which are themselves strings containing
connectives and parentheses \cite[ch.~A.1, \S4, pp.~23, 28]{Barwise1977FirstOrder}.

\emph{Purpose.}  This section states the assumption under which Barwise's
recursion is sound, proves readability under it, shows by an example that
it can fail without the assumption, and proves readability in Barwise's own
application to first-order formulas.  It then derives the recursion theorem
that makes \(\bar v\) a function, the coincidence property behind truth
tables, and the correctness of the truth-table test for tautologies and
consistency that Barwise asserts without proof
\cite[ch.~A.1, \S4, p.~25]{Barwise1977FirstOrder}.
\emph{Inputs.}  Barwise's definitions of propositional formula, truth
assignment, tautology and consistency
(\cref{def:p02-w02-p02-005-bw-propformula,def:p02-w02-p02-005-bw-truth-assignment,def:p02-w02-p02-005-bw-tautology})
and, for the application, his first-order syntax
(\cref{def:p02-w02-p02-005-bw-syntax}).
\emph{Conclusions.}
\Cref{thm:p02-w02b-p02-001-letter-readability,thm:p02-w02b-p02-001-fo-skeleton,thm:p02-w02b-p02-001-recursion,thm:p02-w02b-p02-001-coincidence,thm:p02-w02b-p02-001-truth-table}.
Throughout, \(\circ\) ranges over the binary connectives
\(\land,\lor,\to\), and the length of an expression is its number of
entries.

\subsection{Prime formulas as single symbols}
\label{sec:p02-w02b-p02-001-letters}

Barwise calls a propositional formula of \(P\) an \emph{expression} and
takes the propositional formulas to be the smallest set of expressions that
contains the members of \(P\) and is closed under forming \(\neg A\),
\((A\land B)\), \((A\lor B)\) and \((A\to B)\)
\cite[ch.~A.1, \S4, p.~23]{Barwise1977FirstOrder}.  When the members of
\(P\) are arbitrary objects, the natural reading treats each of them as one
symbol.

\begin{hbconvention}{def:p02-w02b-p02-001-letter-reading}{Expressions over a set of prime formulas}
Let \(P\) be a set with no member among the six symbols
\(\neg,\land,\lor,\to,(,)\).  An \emph{expression of \(P\)} is a finite
sequence each of whose entries is a member of \(P\) or one of these six
symbols.  A member \(p\in P\) is identified with the expression of length
\(1\) whose only entry is \(p\), and \(\neg A\), \((A\circ B)\) denote
concatenations.  The propositional formulas of \(P\) are then those of
\cref{def:p02-w02-p02-005-bw-propformula}.  No restriction is placed on the
cardinality of \(P\).
\end{hbconvention}

\begin{hbtheorem}{thm:p02-w02b-p02-001-letter-readability}{proved=annex:p02-w02b-p02-001-readability}{Unique readability over an arbitrary set of prime formulas}
Under \cref{def:p02-w02b-p02-001-letter-reading}, every propositional
formula \(A\) of \(P\) is exactly one of the following: a member of \(P\);
\(\neg B\) for a propositional formula \(B\); \((B\circ D)\) for a binary
connective \(\circ\) and propositional formulas \(B,D\).  In the second case
\(B\) is uniquely determined by \(A\), and in the third case \(\circ\),
\(B\) and \(D\) are.  Moreover no propositional formula is a proper initial
segment of another.
\end{hbtheorem}

The proof in the annex compares two formulas symbol by symbol from the
left, using induction on the length of the shorter one; it does not use the
parenthesis count of \cref{lem:paren-count}, and it needs no enumeration of
\(P\), so it applies to uncountable sets of prime formulas, the case used in
the second proof of compactness and in the reduction of first-order logic
(\cref{thm:p02-w02-p02-005-bw-prop-compactness,thm:p02-w02-p02-005-bw-main-lemma}).

\begin{hbwarning}{def:p02-w02b-p02-001-string-primes}{Compound strings as prime formulas}
The disjointness assumption cannot simply be dropped when the members of
\(P\) are strings over the same symbols.  Let \(p,q\) be distinct letters
and let \(P=\{p,q,(p\land q)\}\), where \((p\land q)\) is the string of
five symbols.  If expressions are formed by concatenating strings, the
string \((p\land q)\) is a propositional formula of \(P\) in two ways: as a
member of \(P\), and as the conjunction of the members \(p\) and \(q\).
Take the truth assignment with \(v(p)=\mathbf f\), \(v(q)=\mathbf t\) and
\(v((p\land q))=\mathbf t\).  The clause for prime formulas demands
\(\bar v((p\land q))=\mathbf t\); the clause for \(\land\) demands
\(\bar v((p\land q))=\mathbf f\), because \(\bar v(p)=\mathbf f\).  So no
function satisfies all of Barwise's clauses for this \(v\), and the
recursion ``on length of formulas'' has no solution.  Readability is
therefore a hypothesis on \(P\), not a consequence of the definition.
Barwise's remark that prime formulas might be sentences of a natural
language \cite[ch.~A.1, \S4, p.~23]{Barwise1977FirstOrder} is covered by
\cref{def:p02-w02b-p02-001-letter-reading} when each sentence is treated as
one symbol; the case where the prime formulas are first-order formulas is
settled by \cref{thm:p02-w02b-p02-001-fo-skeleton}.
\end{hbwarning}

\subsection{Barwise's application: first-order formulas as propositional formulas}
\label{sec:p02-w02b-p02-001-skeleton}

In the application, \(P\) is the set of formulas of a first-order language
\(L\) that are atomic or begin with a quantifier, and a \emph{tautology of
first-order logic} is a formula of \(L\) that is true under every truth
assignment to this \(P\) \cite[ch.~A.1, \S4, p.~28]{Barwise1977FirstOrder}.
In Barwise's official notation the quantified formulas are
\((\exists v\varphi)\) and \((\forall v\varphi)\), with an outer pair of
parentheses, an equation is \((t_1=t_2)\), and a term built by an \(n\)-ary
function symbol is \(f(t_1\cdots t_n)\), without separating commas
\cite[ch.~A.1, \S3, Definitions~3.2--3.4, pp.~18--19]{Barwise1977FirstOrder}.
So three kinds of formula begin with a left parenthesis, and a prime
formula is a string that contains connectives.  The next theorem shows
that the propositional structure is nevertheless unambiguous.

Barwise lets \(L\) be a set of function, relation and constant symbols,
each function and relation symbol carrying a positive integer arity
\cite[ch.~A.1, \S3, p.~17]{Barwise1977FirstOrder}.  The theorem assumes in
addition what his formation rules tacitly require and the chapter does not
state: the logical symbols \(\land,\lor,\neg,\to,=,\forall,\exists\), the
variables, the two parentheses and the members of \(L\) are pairwise
distinct symbols, and no member of \(L\) is of two kinds.

\begin{hbtheorem}{thm:p02-w02b-p02-001-fo-skeleton}{proved=annex:p02-w02b-p02-001-fo-read}{Readability of first-order formulas over Barwise's prime formulas}
Let \(L\) be a language as just described, with terms, atomic formulas and
formulas as in \cref{def:p02-w02-p02-005-bw-syntax}, and let \(P\) be the
set of formulas of \(L\) that are atomic or of the form \((\exists v\varphi)\)
or \((\forall v\varphi)\).  Form propositional formulas of \(P\) as in
\cref{def:p02-w02-p02-005-bw-propformula}, with concatenation of strings of
symbols.
\begin{enumerate}
\item The propositional formulas of \(P\) are exactly the formulas of \(L\).
\item Every formula of \(L\) is exactly one of: a member of \(P\);
      \(\neg\varphi\) for a formula \(\varphi\); \((\varphi\circ\psi)\) for a
      binary connective \(\circ\) and formulas \(\varphi,\psi\).  In the
      second case \(\varphi\) is uniquely determined, and in the third case
      \(\circ\), \(\varphi\) and \(\psi\) are.
\item No term of \(L\) is a proper initial segment of another term, and no
      formula of \(L\) is a proper initial segment of another formula.
\end{enumerate}
\end{hbtheorem}

Clause~2 is the property that the warning
\cref{def:p02-w02b-p02-001-string-primes} shows can fail for strings in
general.  Here it holds because the second symbol separates the three
kinds of formula that begin with a left parenthesis: after the left
parenthesis an equation continues with the first symbol of a term (a
variable, a constant symbol or a function symbol), a quantified formula
with \(\exists\) or \(\forall\), and a binary compound with the first symbol
of a formula (a left parenthesis, a relation symbol or \(\neg\)).

\subsection{The recursion, prime constituents and truth tables}
\label{sec:p02-w02b-p02-001-recursion}

Say that a set of propositional formulas of \(P\) \emph{reads uniquely} if
each member is exactly one of a member of \(P\), \(\neg B\), or \((B\circ D)\),
with the components uniquely determined.  By
\cref{thm:p02-w02b-p02-001-letter-readability,thm:p02-w02b-p02-001-fo-skeleton},
this holds under \cref{def:p02-w02b-p02-001-letter-reading} and in Barwise's
application.

\begin{hbtheorem}{thm:p02-w02b-p02-001-recursion}{proved=annex:p02-w02b-p02-001-recursion}{Recursion on propositional formulas}
Suppose the propositional formulas of \(P\) read uniquely.  Let \(Y\) be a
set and let \(g_0:P\to Y\), \(g_\neg:Y\to Y\) and
\(g_\circ:Y\times Y\to Y\) for each binary connective \(\circ\) be
functions.  There is exactly one function \(h\) on the propositional
formulas of \(P\) such that \(h(p)=g_0(p)\) for \(p\in P\),
\(h(\neg B)=g_\neg(h(B))\), and \(h((B\circ D))=g_\circ(h(B),h(D))\).
In particular each truth assignment \(v:P\to\{\mathbf t,\mathbf f\}\) has
exactly one extension \(\bar v\) satisfying Barwise's clauses of
\cref{def:p02-w02-p02-005-bw-truth-assignment}.
\end{hbtheorem}

\begin{hbdefinition}{def:p02-w02b-p02-001-constituents}{Prime constituents, recursively}
Barwise describes the prime constituents of \(A\) as the prime formulas out
of which \(A\) is built \cite[ch.~A.1, \S4, p.~23]{Barwise1977FirstOrder}.
When the propositional formulas of \(P\) read uniquely,
\cref{thm:p02-w02b-p02-001-recursion} with \(Y\) the set of finite subsets
of \(P\) makes this precise: the prime constituents of \(p\in P\) form the
set \(\{p\}\); those of \(\neg B\) are those of \(B\); those of \((B\circ D)\)
are those of \(B\) together with those of \(D\).  Every formula has finitely
many prime constituents.
\end{hbdefinition}

\begin{hblemma}{thm:p02-w02b-p02-001-coincidence}{proved=annex:p02-w02b-p02-001-coincidence}{Coincidence on prime constituents}
Suppose the propositional formulas of \(P\) read uniquely.  If two truth
assignments \(v,w:P\to\{\mathbf t,\mathbf f\}\) agree on the prime
constituents of \(A\), then \(\bar v(A)=\bar w(A)\).
\end{hblemma}

This is Barwise's statement that truth tables show ``how the truth or
falsity of a propositional formula depends on the truth or falsity of its
prime constituents'' \cite[ch.~A.1, \S4, p.~24]{Barwise1977FirstOrder}, in
the form needed when \(P\) is infinite; for the handbook's countable
alphabet it is \cref{lem:prop-coincidence}.

\begin{hbtheorem}{thm:p02-w02b-p02-001-truth-table}{proved=annex:p02-w02b-p02-001-tables}{The truth-table test}
Suppose the propositional formulas of \(P\) read uniquely, and let the prime
constituents of \(A\) be the distinct \(q_1,\ldots,q_k\).  Call a function
\(r:\{q_1,\ldots,q_k\}\to\{\mathbf t,\mathbf f\}\) a \emph{row}, and let the
value of \(A\) at \(r\) be \(\bar v(A)\) for any truth assignment \(v\)
extending \(r\).  This value is well defined; there are \(2^k\) rows; \(A\)
is a tautology in the sense of \cref{def:p02-w02-p02-005-bw-tautology} if
and only if its value at every row is \(\mathbf t\); and \(A\) is
consistent if and only if its value at some row is \(\mathbf t\).
\end{hbtheorem}

The theorem justifies Barwise's claim that the method of truth tables makes
it ``a trivial matter'' to decide whether a propositional formula is a
tautology or consistent \cite[ch.~A.1, \S4, p.~25]{Barwise1977FirstOrder}:
the test inspects \(2^k\) rows, and computes each row's value from the
clauses of \cref{def:p02-w02-p02-005-bw-truth-assignment} applied to the
subformulas, which is what the columns of his tables record
(\cref{def:p02-w02-p02-005-bw-table-two,def:p02-w02-p02-005-bw-table-three}).
It decides tautologyhood of a single formula; it does not decide whether an
infinite set of formulas is consistent, which is the question
\cref{thm:p02-w02-p02-005-bw-prop-compactness} reduces to finite subsets.

\begin{hbexample}{def:p02-w02b-p02-001-fo-tautologies}{Prime constituents of first-order formulas}
Barwise gives five sentences to separate the two notions.  He lists
\(\forall x\,R(x)\lor\neg\forall x\,R(x)\) and
\(\neg(\forall x\,R(x)\land\exists x\,S(x))\leftrightarrow
(\neg\forall x\,R(x)\lor\neg\exists x\,S(x))\) as tautologies, and
\((c=c)\), \(\forall x\,(R(x)\lor\neg R(x))\) and
\(\neg\exists x\,S(x)\to\forall x\,\neg S(x)\) as sentences that are not
\cite[ch.~A.1, \S4, p.~28]{Barwise1977FirstOrder}; here
\(A\leftrightarrow B\) abbreviates \((A\to B)\land(B\to A)\)
\cite[ch.~A.1, \S4, p.~25]{Barwise1977FirstOrder}, and the formulas are
written in the informal notation, outer parentheses of quantified formulas
dropped.  By \cref{thm:p02-w02b-p02-001-fo-skeleton}, each has a unique
decomposition into members of \(P\), so
\cref{thm:p02-w02b-p02-001-truth-table} applies.
\begin{enumerate}
\item The first has the single prime constituent \(a=\forall x\,R(x)\) and
      the form \(a\lor\neg a\); both rows give \(\mathbf t\).
\item The second has the two prime constituents \(a=\forall x\,R(x)\) and
      \(b=\exists x\,S(x)\) and the form
      \(\neg(a\land b)\leftrightarrow(\neg a\lor\neg b)\), the first de
      Morgan law of \cref{def:p02-w02-p02-005-bw-tautology}; all four rows
      give \(\mathbf t\).
\item \((c=c)\) is atomic and \(\forall x\,(R(x)\lor\neg R(x))\) begins with a
      quantifier, so each is a member of \(P\) and is its own only prime
      constituent; the row assigning \(\mathbf f\) to it gives
      \(\mathbf f\).  The connectives inside \(\forall x\,(R(x)\lor\neg R(x))\)
      are invisible to the test, because by clause~2 of
      \cref{thm:p02-w02b-p02-001-fo-skeleton} the formula is not of the
      form \(\neg\varphi\) or \((\varphi\circ\psi)\).
\item \(\neg\exists x\,S(x)\to\forall x\,\neg S(x)\) has the two distinct
      prime constituents \(b=\exists x\,S(x)\) and
      \(c'=\forall x\,\neg S(x)\) and the form \(\neg b\to c'\); the row
      \(b\mapsto\mathbf f\), \(c'\mapsto\mathbf f\) gives \(\mathbf f\).
\end{enumerate}
The three non-tautologies are nevertheless true in every structure under
every assignment: a tautology of first-order logic is a much narrower
notion than validity, which is the point Barwise draws from these examples
\cite[ch.~A.1, \S4, p.~28]{Barwise1977FirstOrder}.
\end{hbexample}

\emph{Relationships.}  For the countable alphabet with \(\bot\), the
corresponding facts are \cref{thm:prop-unique-readability} and
\cref{lem:prop-coincidence}; the present results extend them to Barwise's
arbitrary \(P\) and to his strings of first-order symbols.  The handbook's
own first-order grammar, with its readability theorem
\cref{thm:fol-unique-readability}, uses a different notation; clause~3 of
\cref{thm:p02-w02b-p02-001-fo-skeleton} is proved directly for Barwise's
notation.  \Cref{thm:p02-w02b-p02-001-fo-skeleton} is used implicitly
wherever Barwise's reduction applies truth assignments to first-order
formulas, as in \cref{thm:p02-w02-p02-005-bw-lemma-43} and
\cref{thm:p02-w02-p02-005-bw-main-lemma}.
